\section*{5.9 延伸阅读}
\addcontentsline{toc}{section}{5.9 延伸阅读}
\label{sec:5.9}

有关矩阵微分的更多详细信息，以及所需线性代数知识的简要回顾，可参阅 Magnus and Neudecker (2007)。自动微分拥有悠久的历史，我们推荐读者参阅 Griewank and Walther (2003)、Griewank and Walther (2008)、Elliott (2009) 以及其中的参考文献。

在机器学习（以及其他学科）中，我们经常需要计算期望，即需要求解如下形式的积分：

\begin{equation}
\mathbb{E}_{\boldsymbol{x}}[f(\boldsymbol{x})] = \int f(\boldsymbol{x}) p(\boldsymbol{x}) \mathrm{d}\boldsymbol{x}
\tag{5.181}\label{eq:5-181}
\end{equation}

即使 $ p(\boldsymbol{x}) $ 具有便于处理的形式（例如高斯分布），该积分通常也无法解析求解。对 $ f $ 进行泰勒级数展开是寻找近似解的一种方法：假设 $ p(\boldsymbol{x}) = \mathcal{N}(\boldsymbol{\mu}, \boldsymbol{\Sigma}) $ 服从高斯分布，则围绕 $ \boldsymbol{\mu} $ 的一阶泰勒级数展开可以在局部将非线性函数 $ f $ 线性化。对于线性函数，若 $ p(\boldsymbol{x}) $ 服从高斯分布，我们可以精确计算其均值（和协方差）（参见第 6.5 节）。 \emph{扩展卡尔曼滤波（extended Kalman filter）} (Maybeck, 1979) 充分利用了这一性质，用于非线性动力系统（也称为“状态空间模型”）中的在线状态估计。用于逼近式~\eqref{eq:5-181} 中积分的其他确定性方法包括不需要任何梯度的 \emph{无迹变换（unscented transform）} (Julier and Uhlmann, 1997)，或者利用二阶泰勒级数展开（需要黑塞矩阵 / Hessian）在众数周围对 $ p(\boldsymbol{x}) $ 进行局部高斯逼近的 \emph{拉普拉斯近似（Laplace approximation）} (MacKay, 2003; Bishop, 2006; Murphy, 2012)。
